\newif\ifglobal

%%%%%%%%%%%%%%%%%%%%%%%
%%% Compile to view %%%
%%%%%%%%%%%%%%%%%%%%%%%

%%%%%%%%%%%%%%%%%%%%%%%%%%%%%%%%%%%%%%%%%%%%%%%%%%%
%%% Readme:                                     %%%
%%% https://www.overleaf.com/read/bwdjvfjksvjy/ %%%
%%%%%%%%%%%%%%%%%%%%%%%%%%%%%%%%%%%%%%%%%%%%%%%%%%%

% >>> M?: marks a module setting number or important notice

% >>> M4: Set {./relative-path} to external files for this module <<<
\def \modulefiles {../ESP32-H2/10-INTMTRX}
% To add an external file, use:
% (Replace '---' with actual filename)
% '\input{\modulefiles/tables/---.tex}' for tables
% '\includegraphics{\modulefiles/figures/---}' for figures
% '\subfile{\modulefiles/---}' for register files


% >>> M1: This module is in global project? <<<
% 'YES' - uncomment; 'NO' - comment out
\globaltrue

\ifglobal
    \documentclass[main\_\_EN.tex]{subfiles}

\else
    \documentclass[openany, 10pt]{book}

    % >>> M2: In ./00-shared/preamble-trm-module, also find
    %         settings for visibility of todo notes, watermark <<<
    \usepackage{preamble-trm-module}

    % >>> M3: Temporary labels in English,
    %         you can add your own labels there <<<
    \myexternaldocument{./00-shared/temp-labels-en}

\fi %\ifglobal


\setcounter{secnumdepth}{4}


% >>> M5: If this module needs any updates in
% './00-shared/preamble-trm-module.sty'
% (include additional package, etc.), contact your module PM
% or create the file './00-shared/changelog.tex' make the
% necessary updates yourself and document them in this file <<<

\raggedright


\begin{document}


% Sets language into which pre-defined bits of text are translated
\selectlanguage{english}

% Do not delete this block
\ifglobal
% Add the following front matter if
% this module is outside of global project
\else
    \InputIfFileExists{readme}{}{}
    {\let\clearpage\relax \listoftodos}
    {\let\clearpage\relax \listofdonetodos}
    \tableofcontents
    %\listoftables
    %\listoffigures
\fi


%%%%%%%%%%%%%%%%%%%%%%%%%
%%% TRM Module Begins %%%
%%%%%%%%%%%%%%%%%%%%%%%%%

\hypertarget{intmtrx}{}
\chapter{Interrupt Matrix (INTMTX)}
\label{mod:intmtrx}
\section{Overview}

The interrupt matrix embedded in \chipname{} routes one or multiple peripheral interrupt sources to any one of the ESP-RISC-V CPU's peripheral interrupts.

The \chipname{} has 65 peripheral interrupt sources that can be routed to any of the 28 CPU peripheral interrupts using the interrupt matrix.
\begin{tiplisting}
This chapter focuses on how to map peripheral interrupt sources to the CPU interrupts. For information about interrupt configuration, vector, and interrupt handling operations recommended by the ISA, please refer to Chapter \ref{mod:riscvcpu} \textit{\nameref{mod:riscvcpu}}  > Section \ref{sec:riscvcpu-int-ctrl} \textit{\nameref{sec:riscvcpu-int-ctrl}}.
\end{tiplisting}
\vspace{-1em}

\section{Interrupt Terminology in \chipname{}}\label{sec:intmtx-int-terminology}

The following terms related to interrupts are defined in the context of the \textit{\chipname{} Technical Reference Manual} to help readers better understand this document:

\subsection{Interrupt}
An interrupt refers to the event or condition that occurs, causing the CPU to temporarily suspend its current execution and handle a higher-priority task. It is a mechanism that allows the CPU to respond to specific events promptly.

The \textit{\chipname{} Technical Reference Manual} may use the term ``interrupt" in a broader sense to refer to both interrupt signal and interrupt source.

\subsection{Interrupt Signal/Interrupt Source}

Interrupt signal and interrupt source are only defined from different perspectives and mean the same thing.

From the perspective of peripherals, interrupt signals are generated by the peripheral's internal interrupt sources and are sent to the interrupt matrix.

From the perspective of the interrupt matrix, it receives the interrupt signals sent from the peripheral and considers them interrupt sources. The interrupt matrix then outputs CPU peripheral interrupt signals to the CPU.

From the perspective of the CPU, the interrupt signals from the interrupt matrix become sources and are sent to the CPU core together with  the core local interrupt sources.



\subsection{Interrupt Flow in \chipname{}}

Figure \ref{fig:interrupt-term} shows the interrupt flow in \chipname{}.
\begin{figure}[H]
    \centering
    \includegraphics[width=1\textwidth]{\modulefiles/figures/ESP32-H2-interrupt-terms.eps}
    \caption{Interrupt Flow in \chipname{}}
    \label{fig:interrupt-term}
\end{figure}

\section{Features}

The interrupt matrix embedded in \chipname{} has the following features:

\begin{itemize}
    \item 65 peripheral interrupt sources accepted as input
    \item 28 CPU peripheral interrupts generated to the CPU as output
    \item Current interrupt status query of peripheral interrupt sources
    \item Multiple interrupt sources mapping to a single CPU interrupt (i.e., shared interrupts)
\end{itemize}

\section{Architecture}

Figure \ref{fig:interrupt-matrix-structure} shows the structure of the interrupt matrix.

You need to configure the \hyperref[sec:intmtrx-reg-summ]{interrupt matrix registers} to map the peripheral interrupt sources to the CPU interrupts.
The Interrupt Matrix Controller in Figure \ref{fig:interrupt-matrix-structure} manages the mapping and sends the interrupt status of each interrupt source to the interrupt status registers which belong to the interrupt matrix registers.
\begin{figure}[H]
    \centering
    \includegraphics[width=0.8\textwidth]{\modulefiles/figures/ESP32-H2-interrupt.eps}
    \caption{Interrupt Matrix Structure}
    \label{fig:interrupt-matrix-structure}
\end{figure}
\section{Functional Description}

\subsection{Peripheral Interrupt Sources}

The \chipname{} has 65 peripheral interrupt sources in total. Table \ref{tab:peripheral-interrupt-config-status-source} lists all these sources and their mapping registers, as well as the interrupt status registers.


\begin{itemize}
    \item Column ``No.": Peripheral interrupt source number, can be 0 \verb+~+ 64.
    \item Column ``Chapter": in which chapter the interrupt source is described in detail.
    \item Column ``Interrupt Source": Name of the peripheral interrupt source.
    \item Column ``Interrupt Source Mapping Register": Registers used to configure routing of the peripheral interrupt sources to the CPU peripheral interrupts.
    \item Column ``Interrupt Status Register": Registers used to reflect the interrupt source status.
    \begin{itemize}
        \item Column ``Interrupt Status Register - Bit": Bit position in status register, indicating the interrupt status.
        \item Column ``Interrupt Status Register - Name": Name of status registers.
    \end{itemize}
\end{itemize}


\clearpage
\newgeometry{top=2cm,bottom=2cm,left=1.5cm,right=1.5cm}
\begin{landscape}
\input{10-INTMTRX/tables/10-INTMTRX-peri-int-config-status-source__EN}
\end{landscape}
\restoregeometry
\clearpage

\subsection{CPU Interrupts}

The \chipname{} implements its interrupt mechanism using an interrupt controller instead of standard RISC-V ISA. The CPU has 32 interrupts, numbered from 0 to 31. The interrupts numbered 0, 3, 4, and 7 are used by the CPU for core-local interrupts (CLINT), while the remaining 28 interrupts (numbered 1, 2, 5, 6, and 8 \verb+~+ 31) are available for use by the interrupt matrix.

Each CPU interrupt has the following properties:
\begin{itemize}
    \item Priority levels from 1 (lowest) to 15 (highest).
    \item Configurable as level-triggered or edge-triggered.
    \item Lower-priority interrupts maskable by setting interrupt threshold.
\end{itemize}

\begin{tiplisting}
For detailed information about the functions and configuration procedure of CPU interrupts, see Chapter \ref{mod:riscvcpu} \textit{\nameref{mod:riscvcpu}} > Section \ref{sec:riscvcpu-int-ctrl}\textit{\nameref{sec:riscvcpu-int-ctrl}}. The configuration registers of CPU interrupts are listed in Section \ref{sec:intpri-reg-summ} \textit{\nameref{sec:intpri-reg-summ}}.
\end{tiplisting}
\vspace{-1em}

\subsection{Assign Peripheral Interrupt Source to CPU Interrupt}

In this section, the following terms are used to describe the operation of the interrupt matrix.

\begin{itemize}

\item \regindex{SOURCE}: stands for a peripheral interrupt source in Table \ref{tab:peripheral-interrupt-config-status-source}.

\item \hyperref[INTMTXCORE0SOURCEINTRMAPREG]{INTMTX\_CORE0\_\regindex{SOURCE}\_INTR\_MAP\_REG}: stands for the interrupt source mapping register for a \regindex{SOURCE}.

\item Num\_P: stands for the index of CPU interrupts which can be 1, 2, 5, 6, 8 \verb+~+ 31.
\item Interrupt\_P: stands for the CPU interrupt numbered as Num\_P.
\end{itemize}

\subsubsection{Assign One SOURCE to CPU Interrupt}

Setting the corresponding source mapping register \hyperref[INTMTXCORE0SOURCEINTRMAPREG]{INTMTX\_CORE0\_\regindex{SOURCE}\_INTR\_MAP\_REG} of \regindex{SOURCE} to Num\_P assigns this interrupt source to Interrupt\_P.

\subsubsection{Assign Multiple SOURCEs to CPU Interrupt}

Setting the corresponding source mapping register INTMTX\_CORE0\_\regindex{SOURCE}\_INTR\_MAP\_REG of each \regindex{SOURCE} to the same Num\_P assigns multiple sources to the same Interrupt\_P. Any of these sources can trigger CPU Interrupt\_P. When an interrupt signal is generated, CPU should check the interrupt status registers to figure out which peripheral generated the interrupt. For more information, see Chapter \ref{mod:riscvcpu} \textit{\nameref{mod:riscvcpu}} > Section \ref{sec:riscvcpu-int-ctrl}\textit{\nameref{sec:riscvcpu-int-ctrl}}.

\subsubsection{Unassign SOURCE}

Writing 0 to the INTMTX\_CORE0\_\regindex{SOURCE}\_INTR\_MAP\_REG register disables the corresponding interrupt source.

\subsection{Query Current Interrupt Status of SOURCE}

After enabling a \regindex{SOURCE}, you can query its current interrupt status by reading the corresponding bit value in \hyperref[regdesc:INTMTXCORE0INTSTATUS0REG]{INTMTX\_CORE0\_INT\_STATUS\_\regindex{n}\_REG} (read only). For the mapping between \hyperref[regdesc:INTMTXCORE0INTSTATUS0REG]{INTMTX\_CORE0\_INT\_STATUS\_\regindex{n}\_REG} and the \regindex{SOURCE}, please refer to Table \ref{tab:peripheral-interrupt-config-status-source}.


% >>> If your module has no registers, delete sample text
%       for Sections Register Summary and Registers below <<<

%%%%%%%%%%%%%%%%%%%%%%%%
%%% Register Summary %%%
%%%%%%%%%%%%%%%%%%%%%%%%

% >>> Update the sample text below and insert
%     the info on register summary into the subfile <<<

\clearpage
\hypertarget{intmtrx-reg-summ}{}
\begin{landscape}
\section{Register Summary}

\subsection{Interrupt Matrix Register Summary}
\label{sec:intmtrx-reg-summ}
The addresses in this section are relative to the interrupt matrix base address provided in Table \ref{tab:sysmem-base-address} in Chapter \ref{mod:sysmem} \textit{\nameref{mod:sysmem}}.

The abbreviations given in Column \textbf{Access} are explained in Section \hyperref[glossary-access-types]{\textit{Access Types for Registers}}.

\subfile{\modulefiles/intmtx-summary__EN}

\subsection{Interrupt Priority Register Summary}
\label{sec:intpri-reg-summ}
The addresses in this section are relative to the interrupt priority base address provided in Table \ref{tab:sysmem-base-address} in Chapter \ref{mod:sysmem} \textit{\nameref{mod:sysmem}}.

The abbreviations given in Column \textbf{Access} are explained in Section \hyperref[glossary-access-types]{\textit{Access Types for Registers}}.

\subfile{\modulefiles/10-INTPRI-reg-summ__EN}

\end{landscape}

%%%%%%%%%%%%%%%%%
%%% Registers %%%
%%%%%%%%%%%%%%%%%

% >>> Update the sample text below and insert
%      the info on registers into the subfile <<<

\clearpage
\section{Registers}

\subsection{Interrupt Matrix Registers}

The addresses in this section are relative to  the interrupt matrix base address provided in Table \ref{tab:sysmem-base-address} in Chapter \ref{mod:sysmem} \textit{\nameref{mod:sysmem}}.

\subfile{\modulefiles/intmtx-reg__EN}

\subsection{Interrupt Priority Registers}
\label{sec:intpri-reg}

The addresses in this section are relative to  the interrupt priority base address provided in Table \ref{tab:sysmem-base-address} in Chapter \ref{mod:sysmem} \textit{\nameref{mod:sysmem}}.

\subfile{\modulefiles/10-INTPRI-reg__EN}


\end{document}
